% 仅在竞赛过程中实际使用 AI 工具时填写、编译并提交。
% 正式文件名必须为：AI工具使用详情.pdf
\documentclass[UTF8,a4paper,11pt,fontset=none]{ctexart}
\usepackage{fontspec}
\usepackage[left=2.5cm,right=2.5cm,top=2.5cm,bottom=2.5cm]{geometry}
\usepackage{booktabs,tabularx,longtable,array,enumitem,hyperref}
\IfFontExistsTF{Noto Serif CJK SC}{%
  \setCJKmainfont{Noto Serif CJK SC}[BoldFont=Noto Serif CJK SC Bold]
  \setCJKsansfont{Noto Sans CJK SC}[BoldFont=Noto Sans CJK SC Bold]
}{%
  \setCJKmainfont{FandolSong-Regular.otf}[BoldFont=FandolSong-Bold.otf]
  \setCJKsansfont{FandolHei-Regular.otf}[BoldFont=FandolHei-Bold.otf]
}
\hypersetup{hidelinks,pdfauthor={},pdftitle={AI工具使用详情}}
\setlength{\parindent}{2em}
\setlength{\parskip}{0pt}
\linespread{1.25}
\renewcommand{\arraystretch}{1.25}
\newcolumntype{Y}{>{\raggedright\arraybackslash}X}

\title{AI工具使用详情}
\author{}
\date{}

\begin{document}
\maketitle

\section{所用 AI 工具}
\begin{longtable}{p{0.23\textwidth}p{0.67\textwidth}}
\toprule
项目 & 填写内容 \\
\midrule
工具名称 & 【填写实际工具名称】 \\
版本或型号 & 【填写实际版本、模型名称或可识别型号】 \\
使用日期 & 【YYYY-MM-DD 至 YYYY-MM-DD】 \\
\bottomrule
\end{longtable}

\section{具体使用目的和环节}
【逐项写明在哪些环节使用，例如语言润色、代码调试、信息整理等。不要只写“辅助论文”。】

\section{主要提示方式与使用过程}
【概括主要提示方式、迭代过程；需要时附典型交互示例。不得出现参赛者姓名、学校、赛区等身份信息。】

\section{采纳、人工修改与核验}
【除单纯语言润色外，逐项说明 AI 输出是否采纳、进行了哪些人工修改，以及如何用数据、程序、推导或实验进行核验。】

\section{最终核验清单}
\begin{enumerate}[label=\arabic*)]
  \item 公式推导、单位、边界条件和符号已人工核对；
  \item 所有实际使用的程序已在独立环境中运行并复现论文结果；
  \item 引用、数据来源和事实陈述已人工核验；
  \item 核心建模与分析由参赛队主导；
  \item 论文中的 AI 工具使用声明与本文件一致。
\end{enumerate}

\end{document}
